\SourceLeaf{250}%
and \(q\) the angle that the first element of this line makes with the first element of a given shortest line issuing from \(A\). Let \(B\) be a fixed point on this line, along which \(q=0\), and let \(C\) be another fixed point on the surface, for which we shall denote the value of \(q\) simply by \(A\). Suppose that the points \(B,C\) are joined by a shortest line. We shall denote the lengths along it measured from \(B\), as in Article 18, in general by \(s\). As there, let \(\theta\) denote the angle that any element \(\mathrm{d}s\) makes with the element \(\mathrm{d}p\). Finally, let \(\theta^0,\theta'\) be the values of the angle \(\theta\) at the points \(B,C\). We thus have a triangle on the curved surface bounded by shortest lines. Its angles at \(B\) and \(C\), which we shall denote simply by those same letters, are respectively what must be added to \(\theta^0\) to make \(180^0\) and the angle \(\theta'\) itself. Anyone examining our analysis will readily see that we understand all the angles to be expressed as numbers instead of in degrees, taking as the unit the angle \(57^0\,17'\,45''\), which corresponds to an arc equal in length to the radius. So, denoting the circumference of the circle by \(2\pi\), we must set
\[
\theta^0=\pi-B,\qquad\theta'=C
\]
Let us now examine the total curvature of this triangle, which is \(=\int k\mathrm{d}\sigma\), where \(\mathrm{d}\sigma\) denotes an element of the triangle’s surface. Since this element is expressed by \(m\mathrm{d}p\mathord{.}\mathrm{d}q\), we must find the integral \(\int\!\int km\mathrm{d}p\mathord{.}\mathrm{d}q\) over the entire surface of the triangle. Let us begin with integration with respect to \(p\). By \(k=-\frac{1}{m}\mathord{.}\frac{\mathrm{d}\mathrm{d}m}{\mathrm{d}p^2}\), this gives \(\mathrm{d}q\mathord{.}(\mathrm{Const.}-\frac{\mathrm{d}m}{\mathrm{d}p})\) for the total curvature of the area between the lines of the first system corresponding to the values \(q,q+\mathrm{d}q\) of the second variable. Since this curvature must vanish at \(p=0\), the constant introduced by integration must equal the value of \(\frac{\mathrm{d}m}{\mathrm{d}p}\) at \(p=0\), that is, one. We therefore have \(\mathrm{d}q(1-\frac{\mathrm{d}m}{\mathrm{d}p})\), where for \(\frac{\mathrm{d}m}{\mathrm{d}p}\) we must take the value corresponding to the end of this area on the line \(CB\). On this line, however, we have \(\frac{\mathrm{d}m}{\mathrm{d}p}\mathord{.}\mathrm{d}q=-\mathrm{d}\theta\) by the preceding article, so that our expression becomes \(\mathrm{d}q+\mathrm{d}\theta\). Carrying out the second integration from \(q=0\) to \(q=A\), we obtain the total curvature of the triangle \(=A+\theta'-\theta^0=A+B+C-\pi\).

The total curvature equals the area of the part of the sphere’s surface corresponding to the triangle, with a positive or negative sign according as the curved surface on which the triangle lies is concave-concave or concave-convex. The unit of area is the square whose side is one unit (the radius of the sphere), so that the entire surface area of the sphere is \(=4\pi\). Thus the part of the sphere’s surface
